\documentclass{article}

\usepackage{amsmath} % math stuff
\usepackage{amssymb} % math stuff
\usepackage{array} % equations and stuff
\usepackage{bm} % bold math
%\usepackage{booktabs} % extra table rule options
%\usepackage{caption} % suppressed table numbering; incompatible with revtex, and longtable, I think
\usepackage{comment} % comment environment
%\usepackage{enumitem} % customization of enumeration, itemize, and description
\usepackage[T1]{fontenc} % font encoding for special characters, must also use scalable font package
\usepackage[margin=0.8in]{geometry} % paper sizes and margins (but be careful not to mess up pre-defined pages)
\usepackage{graphicx} % for graphics
%\usepackage{helvet} % default font is the helvetica postscript font
\usepackage[utf8]{inputenc} % special characters in tex input
\usepackage{layouts} % print units like widths
\usepackage{lipsum} % lorem ipsum filler text
\usepackage{lmodern} % scalable font?
\usepackage{longtable} % multi-page tables
\usepackage{makecell} % specify line-breaks in table cells
\usepackage{mathrsfs} % math script font
\usepackage{mhchem} % easier chemical formula
\usepackage{microtype} % allows disabling of ligatures
\usepackage{multicol} % multicolumns
%\usepackage{newcent} % new century schoolbook font
\usepackage{nicefrac}
\usepackage{numprint} % print and format (large) numbers
\usepackage{parskip} % removes paragraph indentation, and adjusts paragraph skip, as well as list items
\usepackage{pdfpages} % add pdf files as pages
%\usepackage{setspace} % adjust text spacing and indents
\usepackage{siunitx} % decimal alignment, and degree symbol
\usepackage{subfigure} % divided figures
%\usepackage{tabu} % extra table options
\usepackage{textcomp} % symbols
\usepackage{threeparttablex} % better footnotes with longtable
\usepackage{titling} % title placement
\usepackage{ulem} % strikethrough text
\usepackage{upgreek} % upright Greek letters
%\usepackage{url} % superceded by hyperref
\usepackage{verbatim} % verbatim environment
\usepackage{xcolor} % colors and color boxes
\usepackage{xspace} % commands that don't eat up white space
\usepackage{hyperref} % links and page setup; should always come last

\hypersetup{
 bookmarks=true,
 colorlinks=true,
 citecolor=blue,
 linkcolor=blue,
 urlcolor=blue,
 pdfstartview={XYZ null null 1.0} % default open view is 100%
}

\DisableLigatures[f,t]{encoding = T1} % disable ff, fi, fl, tt ligatures; without options, it also disables -- = endash
\renewcommand{\arraystretch}{1.0} % extra vertical (and horizontal?) space in tables

% define centered, left- and right-aligned columns with specified widths
\newcommand{\PreserveBackslash}[1]{\let\temp=\\#1\let\\=\temp}
\newcolumntype{C}[1]{>{\PreserveBackslash\centering}p{#1}}
\newcolumntype{L}[1]{>{\PreserveBackslash\raggedright}p{#1}}
\newcolumntype{R}[1]{>{\PreserveBackslash\raggedleft}p{#1}}

\begin{document}

\pagestyle{empty} % don't number pages

% custom title
\begin{center}
{\LARGE Classic Riddler}

\vspace{0.15in}

{\Large 18 March 2022}
\end{center}


\section*{Riddle:}

A postal worker and his customer joke about the various ways the customer could mathematically encode her post office box number.

The customer realizes that every integer greater than 1 can be encoded via at least one Fibonacci-like sequence using an ordered triple ($m$, $n$, $q$).
The encoded number is the $q^{\mathrm{th}}$ member of the sequence after the first two positive integers $m$ and $n$, where each term is the sum of the previous two terms.
For example, 7 has the encodings (3, 4, 1) and (1, 3, 2).

In an attempt to stump the postal worker, the customer prefers encodings with a maximal value of $q$.
What encoding should she use for the number 81?

\textit{Extra credit}: What encoding should she use for the number 179?


\section*{Solution:}

For any Fibonacci-like sequence, the ratio of consecutive numbers will always approach the golden ratio, $\phi\approx1.618$.
Therefore, for the number 81 to be as far as possible along in the sequence, its ratio to the previous number should be close to this value.
Specifically, $80/1.618\approx50.1$, so the best candidates for the next-to-last number are 50 and 51.
Working backwards from 50, the reversed sequence would be 81,50,31,19,12,7,5,2,3.
This is as far as it can go with positive integers to start.
This has a $q$ of 7.
For 51, the reversed sequence would be 81,51,30,21,9,12, which has a $q$ of 4.
So the solution to the first question is \fcolorbox{red}{white}{\textbf{(3,2,7)}}\,.

Using the same reasoning for 179, $179/1.618\approx110.6$.
With 110, the reversed sequence would be 179,110,69,41,28,13,15, which has a $q$ of 5.
With 111, the reversed sequence would be 179,111,68,43,25,18,7,11, which has a $q$ of 6.
So the solution to the second question is \fcolorbox{red}{white}{\textbf{(11,7,6)}}\,.

I wrote a table in \texttt{postal\_sums.xlsx} that calculates the maximum $q$ value for combinations of $m$ and $n$ up to 20 and 20.
This confirms the manual solutions above.


\end{document}